\chapter[INTRODUCTION]{Introduction}
\label{ch:introduction}

% Sources for every figure and claim in this chapter were checked against the
% cited papers in October 2026. Counts of reviewed studies (33 corrosion-sensing
% studies, 40 fusion studies) are taken from the literature tables in
% Chapter~\ref{ch:literature-review}.

\section{Background}
\label{sec:background}

Corrosion is the deterioration of a metal through electrochemical interaction
with its environment \citep{Alvani2026}, as the metal returns to a more
stable, oxidized state \citep{Elizalde2023}. It affects
steel structures in every sector that relies on them, and its cost is
substantial. The annual cost of corrosion is usually estimated at 3--4\% of
gross domestic product (GDP) for most nations, of which 15--35\% is thought to
be avoidable, with a significant proportion of the total relating to the cost
of inspection \citep{Nash2022}. Globally the figure has been put at about
US\$2.5 trillion a year, or 3.4\% of world GDP \citep{Imran2025}. Where a
national economy has been surveyed in detail, the estimate holds: the cost to
China in 2014 was approximately 2127.8 billion RMB (about US\$310 billion), or
3.34\% of GDP \citep{Hou2017}. The burden is concentrated in the energy and
maritime sectors. Corrosion accounts for nearly 27\% of gas pipeline
incidents recorded by the European Gas Pipeline Incident Data Group
\citep{Akhlaghi2023}, and corrosion-related expenditure for oil tankers alone
is estimated at US\$969 million a year \citep{Imran2025}. The consequences are
not only financial: pipeline corrosion can lead to leaks, fires, explosions and
the pollution of soil and sea \citep{Akhlaghi2023,Alvani2026}.

Corrosion takes several forms, and they do not present the same inspection
problem. In uniform corrosion, metal is lost evenly across an exposed surface,
so what is visible broadly reflects what has been lost. Pitting is localized:
once the protective film breaks down at a point, attack concentrates there and
forms deep wells of small diameter \citep{Soares2023}.
Pitting is therefore regarded as more dangerous than uniform corrosion,
because it is harder to detect \citep{Soares2023}, and a surface that appears
largely sound can hide deep local loss; in field data from buried transmission
pipelines, pit depths ranged from 0.41 to 13.44\,mm \citep{Akhlaghi2023}. For
this reason the depth of metal loss matters to safe operation as well as the
area it covers \citep{Akhlaghi2023,Huang2026a}, and an inspection method that
reports only the corroded area describes half of the damage.

Inspection is still performed predominantly by manual visual examination,
supported where access permits by non-destructive testing (NDT) such as
ultrasonic, radiographic and magnetic particle inspection \citep{Alvani2026}.
These techniques are effective under controlled conditions but are
time-consuming, costly and dependent on skilled personnel \citep{Alvani2026},
and visual inspection is subjective and inconsistent between inspectors and
conditions \citep{Espinoza2026,Savino2025}. Access compounds the problem.
Inside the storage tanks of a floating production, storage and offloading
(FPSO) vessel, a complete inspection takes ten to fifteen days, because
inspection points must be located precisely in confined spaces of tens of
thousands of cubic meters, and the environment is unhealthy for the inspectors
who enter it \citep{Soares2023}.

Automated visual inspection has consequently received sustained attention.
Deep learning models trained on photographs now detect, segment and grade
corrosion on pipelines, ships, bridges and transmission towers
\citep{Nash2022,Fu2024,Alvani2026,Savino2025}, and an unmanned aerial vehicle
(UAV) can carry the camera to structures that are hazardous or slow to reach
on foot \citep{Savino2025,Nooralishahi2022}. A color camera, however,
records appearance only. Rust can resemble paint, shadows and scratches, and
red, orange and brown coatings can produce color responses very similar to
iron oxide \citep{Espinoza2026,Lavadiya2022}. Corrosion developing beneath an
intact coating presents no visible signature to a camera
\citep{Elizalde2023,Hernndez-Surez2026}. Image quality also degrades in the
environments where corrosion is most severe: salt spray and humidity reduce
contrast and blur corrosion boundaries, and illumination varies
\citep{Mao2026}. Finally, image-based severity is an area measure. Severity is
commonly estimated as the fraction of corroded pixels, which does not measure
material loss or corrosion depth \citep{Alvani2026}.

Other sensors respond to what a camera misses. Infrared thermography responds
to defects beneath the surface through their effect on heat flow rather than
through surface color, and has revealed corrosion hidden under paint and
coatings \citep{Elizalde2023,Lim2021,Doshvarpassand2022}. LiDAR measures
the geometry of a surface at a distance and has been used to assess section
loss in steel members \citep{Maru2025} and to map corrosion in three
dimensions \citep{PimenteldeFigueiredo2025}. No single sensor is sufficient on
its own: RGB cameras are affected by poor lighting and weather, thermal
cameras can be misled by reflective surfaces, and LiDAR point clouds are
sparse, but fusing several sensors allows each to offset the disadvantages of
the others \citep{Tran2024}. This research takes that direction for corrosion
inspection, combining RGB, thermal and LiDAR data in a single detection
framework.

\section{Problem Statement}
\label{sec:problem-statement}

Corrosion is inspected by visual checks and by a range of NDT methods, among
them ultrasonic thickness gauging, guided waves, thermography and hyperspectral
imaging \citep{Elizalde2023,Egodawela2025,Ding2025,Arias2025}. Camera-based
deep learning can locate corrosion on a surface without contact, but it can
mistake coatings, stains and shadows for corrosion
\citep{Espinoza2026,Lavadiya2022}, cannot see beneath intact paint
\citep{Elizalde2023,Hernndez-Surez2026}, and mostly reports the affected area
rather than the depth of loss \citep{Alvani2026,Savino2025}. Depth matters to
safe operation \citep{Akhlaghi2023,Huang2026a}, yet the few studies that
estimate it without contact do so on prepared specimens under controlled
conditions \citep{Egodawela2025,Doshvarpassand2022,Haitz2022,Maru2025}.

Multi-sensor work on corrosion remains limited. Of the 33 corrosion-sensing
studies reviewed in Chapter~\ref{ch:literature-review}, 16 rely on RGB images
alone and five combine more than one sensor. Only one of those five adds
thermal imaging to a camera \citep{Lim2021}, and none fuses RGB, thermal and
LiDAR data. The fusion literature shows the same pattern from the other side:
of the 40 sensor-fusion studies reviewed, three combine RGB, thermal and LiDAR,
and all three apply the combination to three-dimensional object detection for
driving or to sensor calibration rather than to corrosion
\citep{Tran2025,Tran2024,Yang2024}. Where a second sensor has been added for
corrosion, its value is rarely measured in isolation. \citet{Haitz2022} added
laser range and intensity data to color images and found the result only
marginally superior to the color images alone, which led them to ask whether
the additional sensor justifies the complex data fusion it requires.

The problem this research addresses is therefore twofold. No study in this
review has fused RGB, thermal and LiDAR data to detect corrosion and measure
its metal loss, and none has measured what each of the three sensors
contributes to detection.

\subsection*{Hypothesis}

RGB, thermal and LiDAR sensors each fail under different conditions
\citep{Tran2024}. A camera is degraded by shadow, poor lighting and salt spray
\citep{Soares2023,Mao2026,Espinoza2026} and can mistake rust-colored paint
for corrosion \citep{Espinoza2026,Lavadiya2022}; thermal imaging responds to
what lies beneath the surface rather than to its color
\citep{Elizalde2023,Doshvarpassand2022,Hernndez-Surez2026}. Thermal imaging is
affected by sunlight, emissivity and, in passive use, by the weather
\citep{Elizalde2023,Egodawela2025}, and can be misled by reflective surfaces
\citep{Tran2024}; LiDAR measures shape directly rather than inferring it from
heat. LiDAR describes surface geometry only, with no color or temperature and no
view of damage below the surface \citep{Ding2025}, and its sparse point clouds limit
detection \citep{Tran2024}; RGB and thermal imaging supply the surface
appearance and what lies under it. Because no two of the three fail in the
same way, it is hypothesized that a model that fuses RGB, thermal and LiDAR
data will detect corrosion more reliably than a model given any one of them
alone, and that the LiDAR will allow the metal loss of each detected region to
be measured as well as its area. If the ablation in Objective~3 finds the
fused model no better than the best single sensor, the hypothesis is rejected
and that result is reported.

\section{Research Questions}
\label{sec:research-questions}

This research is guided by three questions, one for each objective:

\begin{enumerate}[label=RQ\arabic*.,leftmargin=*,itemsep=8pt]
  \item How can RGB, thermal and LiDAR data of corroded metal surfaces be
        registered into a common reference frame with pixel-level accuracy,
        given that thermal and LiDAR data share little visual appearance with
        the RGB image?
  \item Does a deep learning model that fuses registered RGB, thermal and
        LiDAR data detect and localize corroded regions more accurately than
        the same model given RGB alone, and how accurately can the corroded
        area and metal loss of each detected region be measured from the LiDAR
        data?
  \item What does each sensing modality contribute to detection when
        single-sensor, paired and full configurations, together with a
        blank-channel control, are trained and evaluated under identical
        conditions?
\end{enumerate}

\section{Research Aim and Objectives}
\label{sec:aim-objectives}

The aim of this research is to develop and evaluate a machine learning
framework that detects corroded regions on metal surfaces through the fusion
of RGB, thermal and LiDAR sensor data, measures the corroded area and metal
loss of each region from the LiDAR, and determines the contribution of each
sensing modality to detection. To achieve this aim, the following objectives
are pursued:

\begin{enumerate}[label=\arabic*),leftmargin=*,itemsep=6pt]
  \item To develop a pipeline that spatially aligns RGB, thermal and LiDAR
        data of metal surfaces into a common reference frame.
  \item To develop a fusion-based deep learning model that detects and
        localizes corroded regions in the aligned trimodal data, and to
        measure the corroded area and metal loss of each detected region from
        the LiDAR data.
  \item To evaluate each modality's contribution to detection through an
        ablation comparing single-sensor, RGB--thermal, RGB--LiDAR and full
        RGB--thermal--LiDAR configurations, with a blank-channel control,
        under identical conditions, so that each added channel is isolated
        rather than credited with the residual of the one before it.
\end{enumerate}

Objective~3 is phrased as an evaluation of each sensor's contribution rather
than a demonstration of improvement, so that a finding that a sensor adds
nothing is a valid outcome.

% Task formulation (decided 10 September 2026): object detection, bounding
% boxes, one class "corroded". Depth/metal loss is read from the LiDAR inside
% each detected box, not regressed by the network (decided 9 October 2026).
% The registered dataset is a prerequisite and a contribution, not a formal
% objective, so that failure to obtain field data does not break an objective.

\section{Significance and Contribution of Research}
\label{sec:significance}

This research is expected to contribute the following:

\begin{enumerate}[label=\arabic*),leftmargin=*,itemsep=5pt]
  \item \textbf{A registered RGB--thermal--LiDAR corrosion dataset.}
        Corrosion data are scarce: the lack of a publicly available dataset
        has left models trained on locally produced data that do not
        generalize \citep{Nash2022}. Public corrosion datasets have since
        appeared \citep{Espinoza2026}, but those identified in
        Chapter~\ref{ch:literature-review} are RGB-only, and none pairs the
        three sensors. The dataset will be
        released if the owners of the scanned structures permit it.
  \item \textbf{A contactless framework that detects corrosion and measures
        it.} One fused RGB, thermal and LiDAR model locates each corroded
        region, and the LiDAR measures the corroded area and metal loss of that
        region in physical units, without contact with the surface. Image-based
        severity estimates are a proportion of corroded pixels, which does not
        measure material loss \citep{Alvani2026}; no study in this review both
        detects corrosion with the three sensors and measures it.
  \item \textbf{A controlled ablation of what each sensor adds to
        detection.} Each sensor alone, the paired configurations and the full
        combination are trained and evaluated under identical conditions, with
        a blank-channel control that separates the information a sensor
        contributes from the effect of a larger input. This answers, for a
        contactless sensor set, the question \citet{Haitz2022} left open:
        whether an additional sensor is worth its acquisition and registration
        cost, including the case in which it is not.
\end{enumerate}

For inspection practice, the outcome indicates which sensors are worth
deploying for corrosion detection, and adds a physical measure of metal loss to
the area that camera-based inspection already reports.

\section{Scope of Study}
\label{sec:scope}

The scope of this research is confined within the following boundaries:

\begin{enumerate}[label=\arabic*),leftmargin=*,itemsep=5pt]
  \item Only contactless optical sensing is investigated: RGB imagery, thermal
        infrared imagery and LiDAR. Ultrasonic, eddy current and magnetic flux
        methods are reviewed as context in Chapter~\ref{ch:literature-review}
        but form no part of the proposed system.
  \item Thermal imaging is chosen in preference to hyperspectral imaging.
        Thermal imaging reveals defects beneath the surface
        \citep{Lim2021,Li2024c,Guedes2025}, whereas hyperspectral cameras are
        costly and bulky, which limits their use in the field \citep{Tang2025}.
  \item Detection places bounding boxes around corroded regions under a single
        class. Pixel-level segmentation is not attempted.
  \item Severity is reported as the corroded area and metal loss of each
        detected region, measured from the LiDAR data in physical units.
        Volume and full three-dimensional reconstruction are not attempted.
  \item Pitting is of particular concern, because pits are deep and narrow
        with small diameters \citep{Soares2023}; metal loss is
        therefore measured alongside area. The substrate and service
        environment are not fixed and will follow the data that can be
        acquired.
  \item Real-time adaptive learning, corrosion prevention, and the prediction
        of corrosion rate or remaining life are outside the scope of this work.
  \item All three objectives depend on a registered RGB--thermal--LiDAR
        corrosion dataset that does not exist in public form. Building it is
        therefore an activity of this research, and it is the largest single
        risk to the schedule in Chapter~\ref{ch:methodology}. Acquisition
        begins with controlled corrosion specimens; if access to corroded field
        structures cannot be secured in time, evaluation remains on such
        specimens and the resulting limit on generalization is reported.
\end{enumerate}
